\subsection{实值函数关于滤子的上极限与下极限}

设 \(f\) 是定义在集合 \(\mathbf{X}\) （它由滤子 \(\mathfrak{G}\) 滤子化）上的实值函数。\(\mathfrak{G}\) 关于关系 \(\supset\) 是一个有向集（第一章，§ 6）。对每个 \(\mathbf{M} \in \mathfrak{G}\) 考虑实数 \(\sup_{x \in \mathbf{M}} f(x)\) ：这样便得到函数 \(\mathbf{M} \to \sup_{x \in \mathbf{M}} f(x)\) （它定义在 \(\mathfrak{G}\) 上而取值于 \(\overline{\mathbf{R}}\) ），根据第 4 节命题 7，它是 \(\mathfrak{G}\) 上的递减函数。于是由第 2 节定理 2，它关于有向集 \(\mathfrak{G}\) 有极限。

定义 6. 实值函数 \(\mathbf{M} \to \sup_{x \in \mathbf{M}} f(x)\) 关于有向集 \(\mathfrak{G}\) 的极限称为 \(f\) 关于滤子 \(\mathfrak{G}\) 的上极限，记作 \(\lim \sup_{\mathfrak{G}} f\) ，或 \(\lim \sup_{x, \mathfrak{G}} f(x)\) 。

\(f\) 关于滤子 \(\mathfrak{G}\) 的下极限类似地定义，并记作 \(\liminf_{\mathfrak{G}} f\) 或 \(\liminf_{x, \mathfrak{G}} f(x)\) 。于是我们有

\[
\left\{ \begin{array}{l} \lim \sup_{\mathfrak{G}} f = \lim_{\mathbf {M} \in \mathfrak {G}}(\sup_{x \in \mathbf {M}}f(x)), \\ \lim \inf_{\mathfrak{G}} f = \lim_{\mathbf {M} \in \mathfrak {G}}(\inf_{x \in \mathbf {M}}f(x)). \end{array} \right.\tag{10}
\]

通常在记号中略去滤子 \(\mathfrak{G}\) ，而简单地写 \(\lim \sup f\) 或 \(\lim \sup_x f(x)\) ，或在不会引起混淆时写 \(\lim \sup f(x)\) 。

由公式 (10) 与定理 1，我们得到

\[
\inf _ {x \in \mathbf {X}} f (x) \leqslant \lim \inf _ {\mathfrak {G}} f \leqslant \lim \sup _ {\mathfrak {G}} f \leqslant \sup _ {x \in \mathbf {X}} f (x).\tag{11}
\]

由第 2 节定理 2，我们也可以写

\[
\left\{\lim \sup _ {\mathfrak {G}} f = \inf _ {\mathbf {M} \in \mathfrak {G}} (\sup _ {x \in \mathbf {M}} f (x)), \right.\tag{12}
\]

\[
\left\{\lim \inf _ {\mathfrak {G}} f = \sup _ {\mathbf {M} \in \mathfrak {G}} (\inf _ {x \in \mathbf {M}} f (x)). \right.
\]

又，在公式 (10) 与 (12) 的右端，可以把滤子 \(\mathfrak{G}\) 换成任意一个基 \(\mathfrak{B}\) （\(\mathfrak{G}\) 的基）。

由 (2) 与 (10)，

\[
\lim \inf _ {\mathfrak {G}} f = - \lim \sup _ {\mathfrak {G}} (- f)\tag{13}
\]

因此我们只需考虑上极限。

定理 3. 实值函数 \(f\) 关于滤子 \(\mathfrak{G}\) 的上极限等于 \(f\) 关于 \(\mathfrak{G}\) 的最大聚值。

设 \(b\) 是 \(f\) 关于 \(\mathfrak{G}\) 的一个聚点。对每个 \(\mathbf{M} \in \mathfrak{G}\) ，\(b\) 属于 \(f(\mathbf{M})\) 的闭包，因此 \(b \leqslant \sup_{x \in \mathbf{M}} f(x)\) ，从而由 (12) 得 \(b \leqslant \limsup_{\mathfrak{G}} f = a\) 。

另一方面，设 V 是 a 在 \(\overline{R}\) 中的任意一个开邻域。那么存在 G 中的一个集合 \(M_{0}\) ，使得对每个 \(M \in G\) （它含于 \(M_{0}\) ），都有 \(\sup_{x \in M} f(x) \in V\) ；由于 V 是开集，\(f(M)\) 与 V 相交，因此 a 是 f 关于 G 的一个聚点，证毕。

推论 1. 为使 \(\limsup_{\mathfrak{G}}f = \liminf_{\mathfrak{G}}f\) ，必须且只需 \(f\) 关于滤子 \(\mathfrak{G}\) 有极限，且此时

\[
\lim _ {\mathfrak {G}} f = \lim \sup _ {\mathfrak {G}} f = \lim \inf _ {\mathfrak {G}} f.
\]

因为 \(\overline{\mathbf{R}}\) 是紧的，滤子基 \(f(\mathfrak{G})\) 有极限点当且仅当它只有一个聚点（第一章，§ 9，第 1 节，定理 1 的推论）。

推论 2. 若 \(\mathfrak{H}\) 是比 \(\mathfrak{G}\) 更细的滤子，则我们有

\[
\lim \inf _ {\mathfrak {G}} f \leqslant \lim \inf _ {\mathfrak {H}} f \leqslant \lim \sup _ {\mathfrak {H}} f \leqslant \lim \sup _ {\mathfrak {G}} f.
\]

因为 \(f\) 关于 \(\mathfrak{H}\) 的每个聚点也是 \(f\) 关于 \(\mathfrak{G}\) 的聚点（第一章，§ 7，第 3 节）。

特别地，若 \(\lim_{\mathfrak{H}}f\) 存在，则

\[
\lim \inf _ {\mathfrak {G}} f \leqslant \lim _ {\mathfrak {H}} f \leqslant \lim \sup _ {\mathfrak {G}} f.
\]

推论 3. 设 A 是滤子 G 中的一个集合，\(\mathfrak{G}_{\mathbf{A}}\) 是 G 在 A 上诱导的滤子，\(f_{\mathbf{A}}\) 是 \(f\) 在 A 上的限制，则

\[
\lim \sup _ {\mathfrak {G} _ {\mathbf {A}}} f _ {\mathbf {A}} = \lim \sup _ {\mathfrak {G}} f.
\]

因为滤子基 \(f(\mathfrak{G})\) 的每个聚点也是滤子基 \(f_{\mathrm{A}}(\mathfrak{G}_{\mathrm{A}})\) 的聚点，反之亦然。

由于这个原因，若 \(f\) 只定义在一个子集 \(A\) （含于 \(X\) 且属于 \(\mathfrak{G}\) ）上，则常滥用记号写 \(\limsup_{\mathfrak{G}} f\) 以代替 \(\limsup_{\mathfrak{G}_{\mathbf{A}}} f_{\mathbf{A}}\) 。

命题 11. 设 \(f\) 与 \(g\) 是定义在滤子化集合 \(X\) 上的两个实值函数。则关系 \(f \leqslant g\) 蕴含

\[
\left\{ \begin{array}{l} \lim \sup f \leqslant \lim \sup g, \\ \lim \inf f \leqslant \lim \inf g. \end{array} \right.\tag{14}
\]

这是关系式 (12) 的直接推论。

当 X 是拓扑空间且 \(\mathfrak{G}\) 是 X 的一点 \(a\) 的邻域滤子时，我们写 \(\limsup_{x\to a}f(x)\) {[}相应地，\(\liminf_{x\to a}f(x)]\) 以代替 \(\limsup_{\mathfrak{G}}f\) {[}相应地，\(\liminf_{\mathfrak{G}}f]\) ；显然我们有

\[
\liminf _ {x \to a} f (x) \leqslant f (a) \leqslant \limsup _ {x \to a} f (x).\tag{15}
\]

更一般地，若 X 是拓扑空间 Y 的子空间，且 G 是一点 \(a \in \overline{X}\) 的邻域滤子在 X 上的迹，我们写 \(\limsup_{x \to a, x \in X} f(x)\) {[}相应地，\(\liminf_{x \to a, x \in X} f(x)\) {]}以代替 \(\limsup_{\mathfrak{G}} f\) {[}相应地，\(\liminf_{\mathfrak{G}} f\) {]}；\(\limsup_{x \to a, x \in X} f(x)\) 称为当 x 趋于 a 且保持在 X 中时 \(f(x)\) 的上极限。若 X 是 \(\{a\}\) 的余集，则在这些记号中我们用`` \(x \neq a\) ''代替`` \(x \in X\) ''。

若 \(\mathbf{A}\) 是 \(\mathbf{X}\) 的一个子集且 \(a \in \overline{\mathbf{A}}\) ，则（定理 3 的推论 2）

\[
\liminf _ {x \to a, x \in \mathbf {X}} f (x) \leqslant \liminf _ {x \to a, x \in \mathbf {A}} f (x) \leqslant \limsup _ {x \to a, x \in \mathbf {A}} f (x) \leqslant \limsup _ {x \to a, x \in \mathbf {X}} f (x).
\]

若 V 是 \(a\) 在 Y 中的一个邻域，则有（定理 3 的推论 3）

\[
\lim \sup _ {x \to a, x \in \mathbf {V} \cap \mathbf {X}} f (x) = \lim \sup _ {x \to a, x \in \mathbf {X}} f (x).
\]

因此，拓扑空间中一点处的上极限与下极限概念，和极限概念一样，具有局部特征。

最后，若 \(\mathfrak{G}\) 是 \(\mathbf{N}\) 上的 Fréchet 滤子，则关于 \(\mathfrak{G}\) ，映射 \(n \to u_n\) （它把 \(\mathbf{N}\) 映入 \(\overline{\mathbf{R}}\) ）的上极限（相应地，下极限）记作 \(\limsup_{n \to \infty} u_n\) （相应地，\(\liminf_{n \to \infty} u_n\) ），并称为实数序列 \(u_n\) 的上极限（相应地，下极限）。

于是，关系式 \(\limsup_{n\to \infty}u_n = a\in \mathbb{R}\) 等价于下述论断：对任意给定的 \(\varepsilon >0\) ，存在整数 \(n_0\) ，使得对每个 \(n\geqslant n_0\) 有 \(u_{n}\leqslant a + \varepsilon\) ，且对无穷多个 \(n\) 的值有 \(u_{n}\geqslant a - \varepsilon\) 。当序列上极限的值为 \(+\infty\) 或 \(-\infty\) 时，其定义可以类似地翻译过来。

给定实值函数的序列 \((f_n)\) （定义在集合 \(\mathbf{X}\) 上），我们用 \(\limsup_{n\to \infty}f_n\) （相应地，\(\liminf_{n\to \infty} f_n\) ）表示在任一点 \(x\in \mathbf{X}\) 处取值为 \(\limsup_{n\to \infty}f_n(x)\) {[}相应地，\(\liminf_{n\to \infty} f_n(x)\) {]}的实值函数。由 (10) 与 (12) 我们推出

\[
\left\{\begin{array}{l}\lim \sup _ {n \rightarrow \infty} f _ {n} = \inf _ {n \in \mathbb {N}} (\sup _ {m \geqslant n} f _ {m}) = \lim _ {n \rightarrow \infty} (\sup _ {m \geqslant n} f _ {m}),\\\lim \inf _ {n \rightarrow \infty} f _ {n} = \sup _ {n \in \mathbb {N}} (\inf _ {m \geqslant n} f _ {m}) = \lim _ {n \rightarrow \infty} (\inf _ {m \geqslant n} f _ {m}),\end{array}\right.\tag{16}
\]

其中的极限是在积空间 \(\overline{\mathbf{R}}^{\mathbf{X}}\) 中取的。序列 \((f_n)\) 在 \(\overline{\mathbf{R}}^{\mathbf{X}}\) 中有极限当且仅当 \(\limsup_{n\to\infty}f_n=\liminf_{n\to\infty}f_n\) （定理 3 的推论 1，以及第一章，§ 7，第 6 节，命题 10 的推论 1）。

